\subsection{Literal braid words for the marked geometric Klein four group}

\begin{result}{A.4.36}{Proved}{Two certified rotation paths identify the actual geometric Klein four group in the finite sphere quotient}
\label{thm:sphere10-geometric-klein-words}
On the ten-point configuration
\[
 (0,\infty,1,-1,i,-i,2,-2,1/2,-1/2)
                                                               \tag{GKW.1}
\]
let $a(z)=-z$ and $b(z)=1/z$.  There is an explicitly
specified orientation-preserving marking of this configuration
for which the actual sphere mapping classes $a,b$ are
represented by the literal Artin words
\[
                  A=B\,R_a\,B^{-1},\qquad
                  B_H=B\,R_b\,B^{-1}.
                                                               \tag{GKW.2}
\]
The signed generator lists for $B,R_a,R_b$ are printed in
(GKW.6).  The source words have lengths $90$ and $86$.
Their images in the actual $630$-state quotient of A.4.32
are the specified geometric Klein four group: their
label permutations are respectively
\[
 (2\ 3)(4\ 5)(6\ 7)(8\ 9),\qquad
 (0\ 1)(4\ 5)(6\ 8)(7\ 9),
                                                               \tag{GKW.3}
\]
their squares are one, and they commute.
The statement identifies the geometric source elements
before passing to a finite quotient.  The finite relations
are also verified independently in every one of its $630$
coordinates.  Their actual section coordinates in $E_4$
are retained in the certificate.
\end{result}

\begin{proof}
\par\medskip\noindent\textbf{The actual rotation paths.}\quad
Use the inverse stereographic parametrization
\[
 z=u+iv\longmapsto
   \frac{(2u,2v,u^2+v^2-1)}{u^2+v^2+1}
       \quad\text{on the unit sphere}.
                                                               \tag{GKW.4}
\]
Thus the ten labelled vectors are
\[
\begin{gathered}
 (0,0,-1),\ (0,0,1),\
 (1,0,0),\ (-1,0,0),\
 (0,1,0),\ (0,-1,0),\\
 (4/5,0,3/5),\ (-4/5,0,3/5),\
 (4/5,0,-3/5),\ (-4/5,0,-3/5).
\end{gathered}
\]
Let $\rho_a(\theta)$ be rotation about the positive $Z$
axis by $\theta$, and let $\rho_b(\theta)$ be rotation
about the positive $X$ axis by $\theta$, with
$0\leq\theta\leq\pi$.
In the coordinate (GKW.4), their endpoint maps are exactly
$z\mapsto-z$ and $z\mapsto1/z$.
These are actual ambient paths of homeomorphisms of the
sphere, and not paths selected to satisfy a finite relation.

Choose the rational orthonormal frame
\[
\begin{gathered}
 p=\frac{(2,3,6)}7,\qquad
 e=\frac{(114,-130,27)}{175},\qquad
 f=\frac{(123,90,-86)}{175},\\
 e\times f=p,\qquad
 Z(r)=\frac{e\cdot r+i\,f\cdot r}{1-p\cdot r}.
\end{gathered}
                                                               \tag{GKW.5}
\]
The pole $p$ lies on none of the two sets of rotation
tracks.  For a rotation about a coordinate axis, that
axis coordinate is constant along the orbit.  Here
$p_Z=6/7$ is different from every $Z$ coordinate of
the ten vectors, and $p_X=2/7$ is different from every
one of their $X$ coordinates.  Thus all trajectories
stay in the displayed complex chart.  Moreover
$1-p\cdot r>0$ throughout, by equality in the unit-vector
Cauchy--Schwarz inequality only at $r=p$.

\par\medskip\noindent\textbf{The explicit marking and the three words.}\quad
The real parts of the ten initial chart coordinates are
distinct.  Their increasing label order is
$[4,7,3,9,0,8,5,2,1,6]$.
Begin with the ordered real configuration whose coordinates
are these ten real parts in increasing order, carrying
labels $0,\ldots,9$ in that order.  Apply the positive
permutation braid $B$ below.  Its final labels have just
the displayed order.  Then move each point vertically
to its actual chart coordinate $Z(r_i)$.
These ten vertical paths lie on distinct vertical lines,
so they are disjoint and create no braid crossings.
This composite path specifies the marking completely.
The positive generator convention is the half twist for
which the strand moving from left to right passes below
the other strand; its relative displacement turns
counterclockwise.

In the following lists an entry $j$ means $\sigma_j$,
and an entry $-j$ means $\sigma_j^{-1}$.  Concatenation
of lists means composition in the right-word convention
of A.4.26:
\[
\begin{aligned}
 B={}&[4,3,2,1,7,6,5,4,3,2,5,4,3,\\
      &\qquad 9,8,7,6,5,4,9,8,7,6,8,7,8],\\
 R_a={}&[2,7,3,4,2,3,5,2,4,6,3,9,5,2,4,8,3,2,7,\\
       &\qquad 6,5,4,3,2,1,2,3,8,-9,7,6,4,5,4,3,2,1,6],\\
 R_b={}&[2,3,4,2,-8,3,-6,2,1,2,3,4,5,6,7,4,8,\\
       &\qquad 9,-8,-6,-7,-8,-9,-8,3,2,1,2,3,4,5,6,7,8].
\end{aligned}                                                   \tag{GKW.6}
\]
The first list is obtained by moving the successive
required labels to the left by adjacent positive
crossings.  It has $26$ letters.  The other two lists
are the $38$ and $34$ crossings of the specified rotation
paths in the chart (GKW.5), verified as follows.

\par\medskip\noindent\textbf{An exact rational crossing certificate.}\quad
Fix two initial unit vectors $r_i,r_j$, and put
$d=r_i-r_j$ and $c=r_i\times r_j$.
For either rotation $\rho(\theta)$, the numerator of the
difference of the real chart coordinates is
$e\cdot\rho(\theta)d+f\cdot\rho(\theta)c$.
The numerator of the difference of their imaginary
coordinates is
$f\cdot\rho(\theta)d-e\cdot\rho(\theta)c$.
Both have the same positive denominator
$(1-p\cdot\rho r_i)(1-p\cdot\rho r_j)$.
The cross-product identity proving these formulas uses
$e\times p=-f$ and $f\times p=e$.

Set $t=\tan(\theta/2)$, so $0\leq t<\infty$.
After multiplying the two numerators by $1+t^2$, they
are rational polynomials $Q_{ij}(t)$ and $M_{ij}(t)$
of degree at most two:
\[
\begin{gathered}
 (1+t^2)\rho_a(\theta)(x,y,z)
      =((1-t^2)x-2ty,\ 2tx+(1-t^2)y,\ (1+t^2)z),\\
 (1+t^2)\rho_b(\theta)(x,y,z)
      =((1+t^2)x,\ (1-t^2)y-2tz,\ 2ty+(1-t^2)z),\\
 Q_{ij}=(1+t^2)
          \bigl(e\cdot\rho d+f\cdot\rho c\bigr),\qquad
 M_{ij}=(1+t^2)
          \bigl(f\cdot\rho d-e\cdot\rho c\bigr).
\end{gathered}                                                  \tag{GKW.7}
\]
All coefficients are rational numbers specified by
(GKW.4)--(GKW.5).  Clearing a positive denominator and
dividing by a positive common factor does not change
their signs.

For each of the $45$ pairs and each axis, the certificate
retains the two integral coefficient lists, every
positive root interval of $Q_{ij}$, and the exact
interval bounds for $Q_{ij}'$ and $M_{ij}$ on that
interval.  The isolating endpoints are rational.
Each interval either is a rational root itself or has
opposite nonzero endpoint values of $Q_{ij}$.
Completeness is checked twice: by exact Sturm counts
and by the discriminant and Vieta signs for the
degree-at-most-two polynomial.
All positive roots are simple; neither endpoint
$\theta=0,\pi$ is a crossing.  All event intervals for
one axis are pairwise disjoint.
The interval bounds for $Q_{ij}'$ and $M_{ij}$ have
nonzero constant signs.

At its root, the signed derivative of the horizontal
coordinate difference has the sign of $Q_{ij}'$,
since all removed factors and $dt/d\theta$ are
positive.  The imaginary difference has the sign of
$M_{ij}$.  Hence the signed Artin crossing is exactly
$-\operatorname{sgn}(Q_{ij}'M_{ij})$ in the stated
positive-generator convention.
Sort the disjoint rational intervals in increasing
order.  At each event, the two labels are checked
to be adjacent in the current horizontal order;
their positions give the generator index.
Updating that order gives precisely the two lists
$R_a,R_b$ in (GKW.6).  The final orders are the
permuted initial orders under (GKW.3).
Thus every crossing, its sign and its position is
decided by exact rational arithmetic.

\par\medskip\noindent\textbf{Why these words are the named geometric elements.}\quad
The projection chart only adds the auxiliary pole to
describe an ordinary plane braid.  By isotopy extension,
the projected puncture tracks can be followed by an
ambient plane isotopy fixing that pole.  On the sphere,
compare that isotopy with the specified rotation
$\rho_a(\theta)$ or $\rho_b(\theta)$.
Composing one with the inverse of the other gives an
isotopy fixing all ten initial labelled points at every
time, because the two isotopies have the identical
puncture tracks.  Their endpoint mapping classes on the
ten-point sphere are therefore equal.
Consequently the actual sphere maps represented by
$R_a,R_b$, in the projected configuration, are exactly
$a,b$.

Conjugating by the marking path adds the prefix $B$
and suffix $B^{-1}$; its two vertical portions add no
crossings.  This proves the source equalities (GKW.2).
The endpoint maps $a,b$ are commuting involutions, and
their displayed permutations generate a faithful
$C_2^2$.  Their image is therefore the named geometric
Klein four group in every quotient retaining these
permutations, including the actual quotient of A.4.32.
This conclusion uses the source path comparison,
not a classification of finite complements.

Finally the program
\path{work/ghost_code/check_sphere10_geometric_klein.py}
replays both words using the full $630$-state
multiplication and weights of A.4.32.
It checks both squares and the commutator in every
coordinate, checks the two exact label permutations,
and computes the actual crossing-mod-eight section
coordinates in the fixed $35$-element basis of A.4.31.
Those coordinates satisfy all three section equations
over $\mathbb Z/4$ directly.
The JSON contains all ninety pair-polynomial records,
all seventy-two event certificates, the marking and
source words, and all finite image coordinates.
The text output and matrix archive retain the actual
execution.  This proves the complete marked
identification claimed.
\end{proof}
